% Tabela-síntese da seção 2.3.3 (Aplicações de machine learning em finanças)
\SinteseInicio{Revisões da literatura sobre \acl{ml} em finanças}{tab:sintese_ml_financas}

\SinteseAutor{aranha2023efficacies} &
Antecedentes do \acs{ml} financeiro. &
Revisão de modelos de risco de crédito e previsão de falência. &
Retoma o particionamento recursivo, aplicado à previsão de falências em meados da década de 1980, como antecedente das árvores de decisão em finanças. &
Situa historicamente o uso dos métodos baseados em árvores em finanças. \\
\SinteseLinha

\SinteseAutor{balasubramanian2024survey} &
Evolução da previsão acionária. &
Revisão sistemática da literatura sobre previsão do mercado de ações. &
Relata uma das primeiras experiências com redes neurais na previsão de retornos de ações, no fim da década de 1980, e a incorporação recente de aprendizado profundo, modelos híbridos e sentimento de mercado. &
Contextualiza a evolução do \acs{ml} e seus limites iniciais. \\
\SinteseLinha

\SinteseAutor{manogna2024bibliometric} &
Aprendizado profundo em finanças. &
Análise bibliométrica com citações, cocitação, acoplamento bibliográfico e coocorrência de termos. &
Indica crescimento das aplicações de aprendizado profundo em finanças desde 2017 e quatro grandes áreas de pesquisa, entre elas otimização de portfólio e gestão de risco; aponta como marco pioneiro a aplicação de redes neurais à previsão de retornos de ações no fim da década de 1980. &
Situa a fase recente do aprendizado profundo e dos modelos híbridos na literatura financeira. \\
\SinteseLinha

\SinteseAutor{kumar2025decennaries} &
Redes neurais em finanças. &
Revisão bibliométrica de 3.106 publicações sobre redes neurais artificiais aplicadas às finanças. &
Identifica o estudo pioneiro de particionamento recursivo na previsão de falências entre as 20 referências mais citadas e com maior conectividade do campo. &
Reforça a trajetória histórica do \acs{ml} em finanças e a variedade de problemas tratados. \\

\SinteseFim
